\documentclass[12pt]{article}
\usepackage[a4paper,margin=1in]{geometry}
\usepackage[T1]{fontenc}
\usepackage{lmodern,microtype}
\usepackage{amsmath,amssymb,amsthm,mathtools}
\usepackage{booktabs,array,enumitem,xcolor}
\usepackage[numbers,sort&compress]{natbib}
\usepackage[colorlinks=true,linkcolor=blue!55!black,citecolor=blue!55!black,urlcolor=blue!55!black]{hyperref}
\linespread{1.07}

\newtheorem{theorem}{Theorem}[section]
\newtheorem{proposition}[theorem]{Proposition}
\newtheorem{corollary}[theorem]{Corollary}
\theoremstyle{remark}
\newtheorem{remark}[theorem]{Remark}

\title{Conjugate-Pair Parity and Edge-Quartet Selection\\
\large A Finite Cross-Scale Canonicity Test for the Physical Packet}
\author{Liang Wang}
\date{July 2026}

\begin{document}
\maketitle

\begin{abstract}
The preceding papers identify a genuine four-mode source--observation packet
at $\sigma=0.01$, but its selection was initially phrased as a match to
post hoc physical eigenvalues.  This paper asks whether a simpler dynamical
rule selects the same object.

Because the physical matrices are real, every nonreal simple eigenvalue is
paired with its complex conjugate.  If the outer spectral edge consists of
four nonreal modes separated from the fifth modulus, its real invariant
packet must have even dimension and, in the present geometry, exactly four
dimensions.  We therefore select the four largest-modulus source-observable
modes.

At scales $\sigma=0.04,0.02,0.01$ and on both physical channels, the four
outer modes form two conjugate pairs, all have nonzero source--observation
residue, and the radial gap after the quartet is positive.  The minimum gap
over six cases is $0.05949$.  The length-three RH-185 candidates have zero
two-sided passes at all audited scales, while the finest length-four branch
has twelve.

This supplies a finite selection principle and explains the parity mismatch
of the odd clock branch.  It is not a uniform edge-gap theorem, not a proof
that every scale has a quartet, and not an arithmetic prime-pair or zeta
identification.
\end{abstract}

\section{Why selection matters}

RH-194 matched the four temporal roots to four physical eigenvalues after the
temporal windows had been selected.  This is strong evidence, but a
post hoc match is not a canonical dynamical definition.  A route toward Gate
A needs a rule stated before matching.

Two physical symmetries are available:
\begin{enumerate}[leftmargin=2em]
\item the matrices at the audited levels are real, so complex modes occur in
 conjugate pairs;
\item the temporal packet is intended to describe the outer spectral edge,
 not an arbitrary interior cloud.
\end{enumerate}
Together they suggest an outer-edge quartet rule.

\section{Conjugation parity}

Let $A\in\mathbb R^{n\times n}$.  If $Av=\lambda v$, then
$A\overline v=\overline\lambda\,\overline v$.  Algebraic Riesz projectors
obey
\begin{equation}\label{eq:projector-conjugation}
 P_{\overline\lambda}=\overline{P_\lambda}
\end{equation}
for simple conjugate modes \citep{Kato1995}.

\begin{theorem}[Nonreal parity]\label{thm:parity}
Let $\Lambda$ be a finite isolated spectral set of a real matrix, invariant
under complex conjugation and containing no real eigenvalues.  Then the
algebraic rank of its Riesz projector is even.  In the simple-spectrum case
the modes can be partitioned into conjugate pairs.
\end{theorem}

\begin{proof}
Conjugation is an antilinear bijection from the eigenspace at $\lambda$ to
the eigenspace at $\overline\lambda$.  Since $\Lambda$ contains no real
points, its elements split into disjoint two-element conjugate orbits, with
equal algebraic multiplicities.
\end{proof}

A real packet of odd dimension may contain real modes.  It cannot, however,
be a conjugation-closed packet consisting only of nonreal edge modes.  This
is the relevant obstruction to a length-three approximation of two pairs.

\section{Outer-edge rule}

Order the eigenvalues by modulus:
\begin{equation}\label{eq:modulus-order}
 |\lambda_1|\ge|\lambda_2|\ge\cdots\ge|\lambda_n|.
\end{equation}
Suppose
\begin{equation}\label{eq:edge-gap}
 |\lambda_4|>|\lambda_5|,
\end{equation}
and the first four eigenvalues are nonreal and conjugation closed.

\begin{proposition}[Edge quartet selection]\label{prop:edge}
Under these hypotheses the outer four-mode Riesz set is uniquely defined by
the modulus edge, is conjugation invariant, and has no three-dimensional
conjugation-closed subpacket containing all outer modes.  If the four modes
have nonzero source-observation residues, the associated canonical packet is
four-dimensional and exactly source-observable.
\end{proposition}

The rule is basis-free and does not use the temporal roots.  It can fail if
the edge gap closes, a real eigenvalue enters the edge, or a residue vanishes.

\section{Three-scale physical audit}

We recompute the base spectra at $\sigma=0.04,0.02,0.01$ for both physical
channels.  The top four modes are always two conjugate pairs.  The radial
gap after the quartet is:
\begin{center}
\begin{tabular}{lrrr}
\toprule
scale & side & $|\lambda_4|-|\lambda_5|$ & minimum quartet residue\\
\midrule
$0.04$ & left & $0.06129$ & positive\\
$0.04$ & right & $0.05949$ & positive\\
$0.02$ & left & $0.32537$ & positive\\
$0.02$ & right & $0.32745$ & positive\\
$0.01$ & left & $0.11341$ & positive\\
$0.01$ & right & $0.11766$ & positive\\
\bottomrule
\end{tabular}
\end{center}
The maximum conjugate-pair mismatch in all six cases is below
$2\times10^{-15}$ in the floating computation, as expected from real input
data.

The selected moduli illustrate the branch:
\begin{align}
\sigma=0.04:&\quad 0.6735,0.6735,0.3595,0.3595,\ldots,\\
\sigma=0.02:&\quad 0.7723,0.7723,0.7181,0.7181,\ldots,\\
\sigma=0.01:&\quad 0.7937,0.7937,0.7328,0.7328,\ldots.
\end{align}
The gap is finite at every audited level, although it is not monotone in
$\sigma$.

\section{Parity versus the length-three clock}

The RH-185 candidate lengths are predeclared by the finite clock-rank rule.
The observed gate counts are:
\begin{center}
\begin{tabular}{lrrr}
\toprule
scale & length & windows & two-sided passes\\
\midrule
$0.04$ & 3 & 18 & 0\\
$0.02$ & 3 & 30 & 0\\
$0.01$ & 3 & 40 & 0\\
$0.01$ & 4 & 38 & 12\\
\bottomrule
\end{tabular}
\end{center}

This does not prove that every length-three construction must fail.  It does
show that the audited failures are consistent with a structural mismatch:
the outer physical object is two nonreal conjugate pairs, while a
three-dimensional real conjugation-closed packet cannot contain it.

The length-four success at the finest scale selects the same quartet as the
outer-edge rule.  RH-194 then shows that its roots match the physical modes.

\section{Canonicity and remaining dependence}

The rule “four largest-modulus source-observable modes” is more intrinsic than
“the four modes nearest the temporal roots.”  It uses only the physical
operator, source/observation visibility, and an outer edge.

It is not yet fully canonical in the infinite program.  One must prove:
\begin{enumerate}[leftmargin=2em]
\item a uniform edge gap or a controlled cluster contour;
\item stability of source-observation residues;
\item compatibility of the edge rule under refinement and renormalization;
\item a relation between this finite edge packet and the later cloud ledger.
\end{enumerate}
The present three-scale audit supports, but does not establish, these claims.

\section{What the parity result does not say}

The conjugate-pair theorem is elementary and local.  It does not say that
the edge quartet is a prime pair, that its phases encode prime gaps, or that
its eigenvalues are zeta zeros.  Those interpretations would require new
arithmetic trace identities and are not present here.

Likewise, the failure of the length-three finite windows is not an all-level
no-go theorem.  A different odd packet could include a real mode, use a
non-real coordinate convention, or target a different spectral region.

\section{Updated route}

The ten-paper batch now has a coherent finite branch:
\begin{equation*}
 \begin{gathered}
 \text{outer edge + conjugation parity}
 \longrightarrow\text{source-observable quartet}\\
 \longrightarrow\text{balanced exact packet}
 \longrightarrow\text{temporal determinant/trace ledger}.
 \end{gathered}
\end{equation*}
The next barrier is no longer identifying a local quartet.  It is proving
that the edge rule and its channel weights survive refinement and can be
assembled into an intrinsic object.  That is still Gate A; Gates B--E,
Hilbert--P\'olya, and RH remain untouched.

\bibliographystyle{plainnat}
\bibliography{references}
\end{document}
